% =============================================================================
% 2.4 Modélisation des données
% =============================================================================
\section{Modélisation des données}
\label{sec:ch2-donnees}

\subsection{Modèle conceptuel}

Le modèle conceptuel identifie les entités du domaine, représentées à la
figure~\ref{fig:modele-conceptuel}.

\begin{figure}[htbp]
  \centering
  \resizebox{\linewidth}{!}{% Modèle conceptuel des données
\begin{tikzpicture}[node distance=11mm and 16mm,
                    card/.style={btplabel, fill=white, inner sep=1.2pt}]

  \node[btpboxstrong, text width=2.0cm] (tenant) {ORGANISATION};

  \node[btpbox, text width=2.0cm, below left=12mm and 20mm of tenant] (user) {UTILISATEUR};
  \node[btpbox, text width=2.0cm, below=12mm of tenant] (doc) {DOCUMENT};
  \node[btpbox, text width=2.0cm, below right=12mm and 20mm of tenant] (dossier) {DOSSIER};

  \node[btpboxalt, text width=2.0cm, below=of doc] (chunk) {EXTRAIT};

  \node[btpboxalt, text width=2.0cm, below left=of dossier] (exig) {EXIGENCE};
  \node[btpboxalt, text width=2.0cm, below=of dossier] (plan) {PLAN};
  \node[btpboxalt, text width=2.0cm, below right=of dossier] (export) {EXPORT};

  \node[btpbox, text width=2.0cm, below=of plan] (section) {SECTION};
  \node[btpbox, text width=2.0cm, below=of section] (revue) {REVUE};

  \draw[btparrow] (tenant) -- node[card, pos=0.75] {1..n} (user);
  \draw[btparrow] (tenant) -- node[card, pos=0.75] {1..n} (doc);
  \draw[btparrow] (tenant) -- node[card, pos=0.75] {1..n} (dossier);
  \draw[btparrow] (doc) -- node[card, pos=0.6] {1..n} (chunk);
  \draw[btparrow] (dossier) -- node[card, pos=0.65] {1..n} (exig);
  \draw[btparrow] (dossier) -- node[card, pos=0.6] {1..n} (plan);
  \draw[btparrow] (dossier) -- node[card, pos=0.65] {1..n} (export);
  \draw[btparrow] (plan) -- node[card, pos=0.6] {1..n} (section);
  \draw[btparrow] (section) -- node[card, pos=0.6] {1..n} (revue);

  % relation réflexive parent-enfant
  \draw[btparrowlight] (chunk.west) .. controls +(-1.2,0.4) and +(-1.2,-0.4) .. (chunk.south west)
    node[card, left, xshift=-9mm, yshift=2mm] {parent / enfant};

  % citation d'un extrait par une section
  \draw[btpdashed] (section.west) -- node[card, above, sloped] {cite} (chunk.south east);

\end{tikzpicture}
}
  \caption{Modèle conceptuel des données}
  \label{fig:modele-conceptuel}
\end{figure}

Les entités se regroupent en quatre grappes.

\begin{description}
  \item[Organisation.] L'organisation cliente de la plateforme isole
        utilisateurs, documents, dossiers et conversations. Ce
        cloisonnement n'est pas une option de déploiement~: c'est une
        contrainte de confidentialité.
  \item[Mémoire documentaire.] Un document déposé est découpé en extraits.
        Un extrait peut être parent d'autres extraits. L'identifiant de
        l'extrait est la clé de citation.
  \item[Dossier en cours.] Un dossier de réponse porte des exigences
        extraites, un plan, des sections générées et des exports. Chaque
        section accumule des revues.
  \item[Dialogue.] Une conversation relie des messages à un dossier et,
        éventuellement, à une section.
\end{description}

\subsection{Relations structurantes}

Trois relations méritent d'être soulignées, car elles portent des
invariants du système.

\begin{itemize}
  \item \btpterm{Parent-enfant entre extraits.} La recherche s'exécute sur
        les enfants~; le contexte transmis au modèle est le parent. Un
        enfant sans parent est une anomalie.
  \item \btpterm{Citation d'un extrait par une section.} Une section
        générée référence des identifiants d'extraits. Une citation dont
        l'identifiant n'existe pas est rejetée par la validation
        déterministe.
  \item \btpterm{Revue d'une section.} Toute décision humaine ---
        approbation, correction, relance --- est une entité, pas un
        attribut écrasé. L'historique est appendice, jamais destructif.
\end{itemize}

\subsection{Modèle logique}

Le modèle logique, implanté dans une base relationnelle, traduit ces
entités en tables avec contraintes d'intégrité. Les extraits y sont
stockés en clair, avec leurs métadonnées (pièce, page, rôle, identifiant
parent). Les vecteurs correspondants vivent dans l'index vectoriel, liés
par le même identifiant. Cette dualité --- texte en relationnel, vecteur
en index --- est volontaire~: elle permet de reconstruire un contexte
citable même si l'index est reconstruit, et d'auditer le texte effectivement
transmis au modèle.

Le détail des colonnes et des contraintes figure au
chapitre~\ref{chap:resultats}, où il accompagne les choix d'implémentation.
